\section{ПРОЕКТИРОВАНИЕ И РАЗРАБОТКА СИСТЕМЫ}

\subsection{Архитектура системы}

Разработанная информационная система имеет трёхзвенную архитектуру, состоящую из серверной части (backend), клиентского веб-приложения (frontend) и модуля машинного обучения (analytics-ml). Компоненты развёрнуты в контейнерах Docker и объединены через обратный прокси Nginx.

На рисунке~\ref{arch} представлена архитектура системы. Nginx является единой точкой входа, принимающей HTTP-запросы от клиентов. Запросы к статическим ресурсам и страницам маршрутизируются к контейнеру frontend (SvelteKit), запросы к API --- к контейнеру backend (Rust/Axum). Backend обращается к PostgreSQL для хранения данных, к Redis для кэширования и rate limiting, и к analytics-ml для получения прогнозов машинного обучения. Синхронизация с облачным API Lely выполняется периодически по таймеру внутри backend.

\begin{figure}
  \centering
  \includegraphics[width=0.95\textwidth,height=0.85\textheight,keepaspectratio]{container.png}
  \caption{Архитектура информационной системы}
  \label{arch}
\end{figure}

Архитектурные решения обоснованы несколькими соображениями. Разделение ответственности обеспечивает чёткое разграничение: backend отвечает за бизнес-логику, frontend --- за пользовательский интерфейс, analytics-ml --- за модели машинного обучения. Независимое масштабирование позволяет запускать несколько экземпляров backend за балансировщиком нагрузки без изменения остальных компонентов. Технологическая гетерогенность даёт возможность использовать наиболее подходящий стек для каждого компонента (Rust для серверной части, Svelte для клиента, Python для ML). Изоляция сбоев, обеспечиваемая контейнерной изоляцией Docker, гарантирует, что сбой одного компонента не приводит к отказу всей системы.

\subsection{Модель данных}

База данных PostgreSQL содержит более 35 таблиц, обеспечивающих хранение всех сущностей предметной области. Перечень таблиц с назначением приведён в приложении~\ref{app-db}, полная ER-модель --- в приложении~\ref{app-er}. На рисунке~\ref{er} представлена модель данных системы.

\begin{figure}
  \centering
  \includegraphics[width=0.95\textwidth,height=0.85\textheight,keepaspectratio]{database-model.png}
  \caption{Модель данных информационной системы}
  \label{er}
\end{figure}

Основные таблицы базы данных и их назначение:

\textbf{users} --- пользователи системы. Содержит имя пользователя, хеш пароля (bcrypt), роль (admin/user), флаг обязательной смены пароля и дату создания. Пароли хранятся в виде хешей с использованием алгоритма bcrypt с автоматической генерацией соли, что исключает утечку паролей при компрометации базы данных. Хеш-функция bcrypt является адаптивной: параметр cost (по умолчанию 12) определяет вычислительную сложность и может быть увеличен по мере роста производительности оборудования.

\textbf{animals} --- реестр животных. Содержит уникальный идентификатор, инвентарный номер (life\_number), кличку, пол, дату рождения, породу, местоположение, статус активности и множество других полей. Поддерживается флаг \texttt{active} для деактивации животных без удаления записей, что обеспечивает сохранение истории. Инвентарный номер имеет уникальное ограничение для предотвращения дублирования.

\textbf{milk\_day\_productions} --- суточные надои. Привязаны к животному и дате. Содержат объём молока (кг), жир (\%), белок (\%), лактозу (\%), количество соматических клеток (\acrshort{scc}), энергию (MJ) и жирово-белковое соотношение (\acrshort{fpr}). Композитный индекс (animal\_id, date) обеспечивает быстрый поиск надоев за период.

\textbf{milk\_visits} --- визиты на доение. Содержат информацию о каждом доении: время начала и окончания, объём молока, скорость доения (кг/мин), данные по четвертям вымени (объём, время, процент выдаивания), флаг успешности доения и идентификатор робота.

\textbf{milk\_quality} --- результаты лабораторных анализов молока. Включает количество соматических клеток (\acrshort{scc}), жир, белок, лактозу, мочевину, бактерии.

\textbf{calvings, calves, inseminations, pregnancies, heats, dry\_offs} --- данные о воспроизводстве. Каждая таблица связана с животным через внешний ключ. Поддерживается полная история всех событий воспроизводства, включая дату события, результат и примечания.

\textbf{feed\_types, feed\_groups, feed\_day\_amounts, feed\_visits} --- данные о кормлении. Типы кормов описывают рационы с указанием стоимости и единиц измерения, группы --- привязку кормов к группам животных, дневные объёмы --- плановое и фактическое потребление, визиты --- отдельные посещения кормостанции с временем и объёмом.

\textbf{activities, ruminations, grazing\_data} --- данные фитнес-мониторинга и пастьбы. Активность измеряется в шагах за период, жвачка --- в минутах жевания за период. Данные поступают от датчиков или через интеграцию с Lely.

\textbf{alerts} --- система оповещений. Оповещения создаются автоматически на основе настраиваемых пороговых значений. Каждое оповещение содержит категорию (milk, health, reproduction, feed, equipment), степень важности (high, medium, low), статус (active, acknowledged, resolved), идентификатор животного и текстовое описание.

\textbf{lely\_sync\_state} --- состояние синхронизации с Lely. Хранит дату последней успешной синхронизации для каждого типа данных, что позволяет инкрементально обновлять информацию и минимизировать объём данных при каждой синхронизации.

\textbf{lely\_config} --- конфигурация интеграции с Lely. Хранит URL API, зашифрованные логин и пароль, ключ фермы и настройки синхронизации. Учётные данные шифруются с использованием AES-256-GCM~\cite{AES-GCM}.

Для эффективной работы с временными рядами (надои, активность, жвачка) используется расширение TimescaleDB~\cite{TimescaleDB}. Таблицы с временными рядами преобразованы в гипертаблицы с автоматической партицировкой по времени (по умолчанию --- 7 дней на чанк), что ускоряет запросы за определённый период, так как сканируются только релевантные чанки. Непрерывные агрегаты (continuous aggregates) обеспечивают предвычисление суточных и часовых агрегатов, обновляемых в фоновом режиме.

Индексация включает уникальные ограничения (инвентарный номер животного), композитные индексы (животное + дата) и частичные индексы (только активные животные). Расширение pg\_trgm~\cite{PgTrgm} обеспечивает полнотекстовый поиск по кличкам животных с использованием триграммного сходства.

\subsection{Серверная часть (Backend)}

Серверная часть реализована на языке Rust с использованием фреймворка Axum. Архитектура следует паттерну разделения на слои: HTTP-обработчики (handlers), сервисы бизнес-логики (services), модели данных (models) и промежуточное программное обеспечение (middleware). На рисунке~\ref{backend-arch} представлена архитектура серверной части.

\begin{figure}
  \centering
  \includegraphics[width=0.95\textwidth,height=0.85\textheight,keepaspectratio]{backend-components.png}
  \caption{Архитектура серверной части}
  \label{backend-arch}
\end{figure}

\subsubsection{Маршрутизация и обработка запросов}

Все API-эндпоинты доступны по префиксу \texttt{/api/v1/}. Маршруты сгруппированы по доменным областям: животные, надои, воспроизводство, кормление, фитнес, отчёты, аналитика, настройки и другие. Общее количество обработчиков --- более 25 модулей. Примеры запросов и ответов приведены в приложении~\ref{app-api}.

Документация API генерируется автоматически с помощью библиотеки utoipa~\cite{utoipa} на основе аннотаций в коде и доступна через Swagger UI по адресу \texttt{/api/v1/docs}. Аннотации включают описание эндпоинтов, типы запросов и ответов, коды ошибок. Документация всегда актуальна, так как генерируется из исходного кода.

\subsubsection{Аутентификация и авторизация}

Аутентификация основана на \acrshort{jwt}~\cite{JWT} (JSON Web Token) с алгоритмом HS256. При входе в систему сервер выдаёт два токена: короткоживущий access token (по умолчанию 15 минут) для авторизации запросов, содержащий имя пользователя (sub), роль (role), флаг обязательной смены пароля (must\_change\_password) и уникальный идентификатор токена (jti); и долгоживущий refresh token (по умолчанию 7 дней) для обновления access token без повторного ввода пароля.

Токены хранятся в HttpOnly, SameSite=Strict cookie, что предотвращает доступ к ним из JavaScript и обеспечивает защиту от CSRF-атак~\cite{OWASP-CSRF}. Access token может передаваться также в заголовке \texttt{Authorization: Bearer <token>} для поддержки программных клиентов.

Отзыв токенов реализован через таблицу \texttt{revoked\_tokens}, хранящую JTI отозванных токенов. При каждом запросе проверяется, не был ли токен отозван. Просроченные записи удаляются ежечасно фоновым заданием. Механизм отмены необходим для реализации безопасного выхода из системы: при выходе оба токена (access и refresh) добавляются в список отозванных.

Авторизация реализована через три уровня доступа, соответствующих типам Axum-экстракторов. Уровень \textbf{Claims} допускает любого аутентифицированного пользователя с валидным access token, при этом запрос блокируется, если требуется сменить пароль. Уровень \textbf{ClaimsAllowMustChange} также требует аутентификации, но допускает состояние обязательной смены пароля и используется для эндпоинта смены пароля. Уровень \textbf{AdminGuard} ограничивает доступ пользователями с ролью \texttt{admin} и применяется для управления пользователями, настройками системы и резервным копированием.

Время жизни \acrshort{jwt} настраивается через таблицу \texttt{system\_settings} без перезапуска сервера. Изменение вступает в силу для новых токенов.

\subsubsection{Промежуточное программное обеспечение (Middleware)}

Стек middleware обеспечивает безопасность, наблюдаемость и ограничение нагрузки. Компоненты middleware применяются в порядке, при котором каждый запрос последовательно проходит через все слои. \textbf{CORS} настраивает кросс-доменный доступ~\cite{CORS} с конфигурируемыми источниками, методами и разрешёнными заголовками; в продакшн-конфигурации разрешены только домены приложения. \textbf{CSRF-защита} проверяет заголовок Origin для изменяющих состояние запросов (POST, PUT, DELETE, PATCH) и отклоняет запрос при отсутствии или несовпадении Origin. \textbf{Rate Limiting} ограничивает частоту запросов по IP-адресу (по умолчанию 100 запросов за 60 секунд) с использованием Redis при наличии или в памяти в качестве fallback; поддерживается определение реального IP через заголовки X-Forwarded-For и X-Real-IP при работе за обратным прокси. \textbf{Gzip-сжатие} уменьшает объём передаваемых данных. \textbf{Логирование} обеспечивает трассировку HTTP-запросов и ответов через tracing с поддержкой структурированных логов в JSON-формате. \textbf{Метрики Prometheus} включают счётчики \texttt{http\_requests\_total} и гистограммы \texttt{http\_request\_duration\_seconds} с нормализацией маршрутов. \textbf{Идентификаторы запросов} передаются в заголовке \texttt{X-Request-Id} для трассировки между сервисами. Наконец, \textbf{ограничение размера тела} (10 МБ) предотвращает атаки с чрезмерной нагрузкой.

\subsubsection{Интеграция с Lely}

Модуль интеграции с Lely обеспечивает автоматическую двустороннюю синхронизацию данных с облачным API Lely Horizon. Синхронизация выполняется по настраиваемому таймеру (по умолчанию каждые 5 минут).

Архитектура модуля включает HTTP-клиент с аутентификацией по токену и автоматическим обновлением каждые 58 минут (токен Lely действует 60 минут), а также повторными попытками с экспоненциальной задержкой (2 попытки, базовая задержка 1 секунда). Планировщик синхронизации поддерживает параллельные задачи с мьютексом для предотвращения одновременного запуска и инкрементальную синхронизацию по временным диапазонам (чанками по 7 дней для временных рядов). Маппинг данных использует кэш инвентарных номеров (HashMap<LifeNumber, AnimalId>) для разрешения Life Number Lely во внутренний идентификатор животного; кэш обновляется перед каждой синхронизацией. Upsert-логика выполняет добавление или обновление записей по уникальному ключу (например, animal\_id + date для суточных надоев) через INSERT ON CONFLICT UPDATE. Учётные данные (логин и пароль Lely) шифруются алгоритмом AES-256-GCM~\cite{AES-GCM}, а ключ шифрования задаётся через переменную окружения \texttt{LELY\_ENCRYPTION\_KEY}.

Синхронизируются 23 типа данных: животные, отёлы, осеменения, стельности, охоты, запуски, суточные надои, визиты на доение, качество молока, данные роботов, кормление, активности, жвачка, пастьба, быки-производители, переводы, кровные линии, типы и группы кормов, контакты, местоположения.

Конфигурация интеграции хранится в таблице \texttt{lely\_config} и может быть изменена через веб-интерфейс без перезапуска сервера. При изменении конфигурации планировщик синхронизации перезапускается через CancellationToken.

\subsection{Клиентское веб-приложение (Frontend)}

Клиентское приложение реализовано на SvelteKit 2 с Svelte 5~\cite{Svelte5} и представляет собой одностраничное приложение (\acrshort{spa}) с серверным рендерингом (\acrshort{ssr}). На рисунке~\ref{frontend-arch} представлена архитектура клиентского приложения.

\begin{figure}
  \centering
  \includegraphics[width=0.95\textwidth,height=0.85\textheight,keepaspectratio]{frontend-structure.png}
  \caption{Архитектура клиентского веб-приложения}
  \label{frontend-arch}
\end{figure}

\subsubsection{Маршрутизация и навигация}

Приложение содержит более 24 страниц, сгруппированных по функциональным областям: дашборд, животные, надои, воспроизводство, кормление, фитнес, пастьба, отчёты, аналитика, настройки, контакты, местоположения, быки, переводы, запасы, оповещения, задачи, ветеринария, финансы, рентабельность, поиск.

Навигация реализована через боковую панель (Sidebar) с иконками из библиотеки Lucide и группировкой разделов. Адаптивная верстка на основе Tailwind CSS обеспечивает корректное отображение на мобильных устройствах: на экранах менее 768px боковая панель сворачивается в гамбургер-меню.

\subsubsection{Аутентификация на стороне клиента}

Серверный хук (hooks.server.ts) проверяет наличие \acrshort{jwt} в cookie при каждом запросе. Декодируется полезная нагрузка токена (без проверки подписи, только чтение полей) и проверяется срок действия. Неаутентифицированные пользователи перенаправляются на страницу входа, аутентифицированные --- не могут перейти на страницу входа.

Управление состоянием аутентификации реализовано через Svelte-стор (auth store), который инициализируется при загрузке приложения через endpoint refresh. Стор содержит имя пользователя, роль, флаг обязательной смены пароля и состояние аутентификации.

\subsubsection{API-клиент}

Клиент для взаимодействия с серверным API реализует автоматическое обновление токена: при получении ответа 401 выполняется попытка обновления access token через \texttt{/auth/refresh}; если refresh-токен также недействителен, пользователь разлогинивается и перенаправляется на страницу входа.

Для повышения надёжности реализован механизм повторных попыток: до 3 попыток с экспоненциальной задержкой (500 мс базовая задержка + случайный джиттер от 0 до 500 мс) для ответов с кодом 5xx. Джиттер предотвращает синхронизацию повторных попыток от нескольких клиентов.

API-клиент разделён на 23 модуля по доменным областям, каждый из которых предоставляет типизированные методы для выполнения CRUD-операций.

\subsubsection{Компоненты пользовательского интерфейса}

Библиотека переиспользуемых компонентов включает DataTable --- таблицу данных с сортировкой по столбцам, серверной пагинацией и адаптивным отображением на мобильных устройствах (на маленьких экранах таблица преобразуется в карточки); Modal и ConfirmDialog --- модальные окна для создания, редактирования и подтверждения действий с поддержкой клавиши Escape; FormField --- компонент формы с поддержкой валидации, отображения ошибок и меток; FilterBar --- панель фильтров для таблиц с поддержкой текстового поиска, выпадающих списков и диапазонов дат; Pagination --- компонент пагинации с отображением номера текущей страницы и общего количества; TabBar --- панель вкладок для переключения между разделами одной страницы; Toaster --- уведомления (success/error/warning/info) с раскрываемыми подробностями и автоматическим скрытием; MilkChart, LactationCurveChart, ReportChart, SeasonalChart --- компоненты визуализации данных на основе Chart.js с адаптивными размерами; OfflineIndicator --- индикатор состояния сети, отображающийся при потере соединения.

Валидация форм реализована на основе Svelte 5 Runes --- реактивной системы состояния, позволяющей декларативно описывать правила валидации и отслеживать состояние формы (dirty, errors, touched). Правила валидации включают проверку обязательных полей, минимальной и максимальной длины, формата даты и числовых значений.

\subsubsection{Генерация отчётов}

Система поддерживает генерацию отчётов более 17 типов, покрывающих все аспекты управления стадом. Каждый отчёт доступен в формате JSON для отображения в интерфейсе, а также поддерживает экспорт в CSV и PDF.

Ключевые типы отчётов перечислены в таблице~\ref{reports}.

\begin{table}
  \caption{Основные типы отчётов системы}
  \begin{tabular}{|c|l|p{6cm}|}
    \hline
    \textbf{Код} & \textbf{Название} & \textbf{Описание} \\
    \hline
    R12 & Здоровье вымени & Рабочий список коров с отклонениями за 24 часа \\
    \hline
    R13 & Неудачные доения & Визиты с неуспешным доением \\
    \hline
    R16 & Обзор стада & Дневная производительность стада \\
    \hline
    R18 & Остаток корма & Невыданный корм по коровам \\
    \hline
    R20 & Надой по времени & Суточный надой с разбивкой по показателям \\
    \hline
    R23 & Анализ вымени & Расширенный анализ за 14 дней \\
    \hline
    R24 & Активность/жвачка & Здоровье, активность и жвачка \\
    \hline
    R31-34 & Календарь & Ожидаемые отёлы, запуски, охота, стельность \\
    \hline
    R41 & Эффективность & Эффективность корова-робот \\
    \hline
    R52 & Лактация & Кривые лактации по номерам \\
    \hline
    R56 & Робот & Производительность роботов \\
    \hline
    R70 & Корм по типам & Стоимость и объём кормов по типам \\
    \hline
    R72 & Корм на корову & Потребление корма на корову \\
    \hline
  \end{tabular}
  \label{reports}
\end{table}

\subsection{Модуль машинного обучения (Analytics-ML)}

Модуль машинного обучения реализован на Python с использованием FastAPI и XGBoost. Он предоставляет API для прогнозирования и обнаружения аномалий, используемый серверной частью. На рисунке~\ref{ml-arch} представлена архитектура модуля машинного обучения.

\begin{figure}
  \centering
  \includegraphics[width=0.95\textwidth,height=0.85\textheight,keepaspectratio]{analytics-ml.png}
  \caption{Архитектура модуля машинного обучения}
  \label{ml-arch}
\end{figure}

\subsubsection{Модели машинного обучения}

Разработано 11 моделей машинного обучения. В таблице~\ref{mlmodels} приведено их краткое описание.

\begin{table}
  \caption{Модели машинного обучения}
  \small
  \begin{tabularx}{\textwidth}{|p{3.5cm}|l|p{3cm}|X|}
    \hline
    \textbf{Модель} & \textbf{Алгоритм} & \textbf{Тип} & \textbf{Назначение} \\
    \hline
    Прогноз надоев (корова) & XGBoost & Регрессия & Дневной надой с q10/q90 доверительными интервалами \\
    \hline
    Прогноз надоев (стадо) & Prophet & Регрессия & Тренд надоя стада с учётом сезонности \\
    \hline
    Риск мастита & XGBoost & Классификация & Вероятность мастита (высокая/средняя/низкая) \\
    \hline
    Предупреждение кетоза & XGBoost & Классификация & Риск метаболического нарушения \\
    \hline
    Обнаружение охоты & XGBoost & Классификация & Состояние (в охоте/приближается/нет) \\
    \hline
    Риск выбраковки & XGBoost & Регрессия & Дни до выбраковки \\
    \hline
    Оценка BCS & XGBoost & Регрессия & Оценка упитанности (1.0--5.0) \\
    \hline
    Кластеризация стада & KMeans & Кластеризация & Сегментация стада \\
    \hline
    Аномалии оборудования & Isolation Forest & Аномалии & Аномалии работы роботов \\
    \hline
    Рекомендации по корму & XGBoost & Регрессия & Рекомендуемый объём корма \\
    \hline
  \end{tabularx}
  \label{mlmodels}
\end{table}

\textbf{Прогноз надоев на корову.} Модель XGBoost с квантильной регрессией предсказывает дневной надой с доверительными интервалами. Целевая переменная --- объём молока (кг) за сутки. Признаки включают: надои за предыдущие 7, 14, 30 дней; скользящие среднее, стандартное отклонение, минимум, максимум за окна 7 и 14 дней; \acrshort{dim} (дни лактации); номер лактации; \acrshort{fpr}; сезонные признаки (месяц, день недели).

Для построения доверительного интервала обучаются три модели: для медианы (q50), нижней границы (q10) и верхней границы (q90). Объективная функция Quantile Loss (пинбольная функция потерь) для квантиля $\alpha$ определяется как:

$$L_\alpha(y, \hat{y}) = \begin{cases} \alpha \cdot (y - \hat{y}), & \text{если } y \geq \hat{y} \\ (1 - \alpha) \cdot (\hat{y} - y), & \text{если } y < \hat{y} \end{cases}$$

где $y$ --- фактическое значение, $\hat{y}$ --- прогноз.

\textbf{Прогноз надоев на стадо.} Модель Prophet~\cite{Prophet} используется для прогнозирования надоев на уровне стада. Модель декомпозирует временной ряд на три компоненты:

$$y(t) = g(t) + s(t) + h(t) + \varepsilon_t$$

где $g(t)$ --- тренд (кусочно-линейная или логистическая функция), $s(t)$ --- сезонная компонента (ряд Фурье), $h(t)$ --- эффект праздников и особых дней, $\varepsilon_t$ --- остатки. Prophet автоматически обнаруживает точки изменения тренда и устойчив к пропущенным данным и выбросам.

\textbf{Риск мастита.} Модель XGBoost Classifier классифицирует вероятность мастита~\cite{MastitisDetection} (высокая/средняя/низкая) на основе \acrshort{scc}, электрической проводимости молока, \acrshort{fpr} и изменений в надоях. Ключевым признаком является динамика \acrshort{scc}: резкое увеличение (более чем на 50~процентов за сутки) является сильным индикатором воспаления вымени.

\textbf{Предупреждение кетоза.} Модель определяет риск метаболического нарушения через \acrshort{fpr}, \acrshort{dim}, надои и потребление корма. Кетоз характерен для первых 60 дней лактации, когда потребность в энергии превышает потребление. Жирово-белковое соотношение (\acrshort{fpr}) выше 1.4 является индикатором субклинического кетоза.

\textbf{Обнаружение охоты.} Модель классифицирует состояние коровы~\cite{PrecisionLivestock} на основе активности (шаги), жвачки (минуты), надоев и визитов. Коровы в охоте демонстрируют повышенную активность (в 3--5 раз выше базовой), снижение жвачки и увеличение количества визитов на доение.

\textbf{Риск выбраковки.} Модель XGBoost Regressor прогнозирует количество дней до выбраковки на основе возраста, надоев, \acrshort{scc}, номера лактации и индекса здоровья. Целевая переменная вычисляется как разница между текущей датой и датой фактической выбраковки для животных с известным исходом.

\textbf{Оценка упитанности (BCS).} Модель оценивает \acrshort{bcs} (1.0--5.0) на основе веса, надоев, \acrshort{dim}, жвачки и эффективности кормления (надой/потребление корма). BCS --- важный показатель здоровья, однако его ручная оценка трудоёмка и субъективна, что делает автоматизированную оценку ценной для фермера.

\textbf{Кластеризация стада.} Алгоритм KMeans~\cite{KMeans} сегментирует стадо на группы с похожими характеристиками. Число кластеров $k$ выбирается автоматически по метрике силуэта:

$$s(i) = \frac{b(i) - a(i)}{\max\{a(i), b(i)\}}$$

где $a(i)$ --- среднее расстояние от объекта $i$ до объектов того же кластера, $b(i)$ --- среднее расстояние от объекта $i$ до объектов ближайшего другого кластера. Значение силуэта находится в диапазоне $[-1, 1]$, при этом значение ближе к 1 указывает на правильную кластеризацию.

\textbf{Аномалии оборудования.} Алгоритм Isolation Forest~\cite{IsolationForest} обнаруживает аномалии в работе доильных роботов (скорость доения, время по четвертям, частота отказов). Алгоритм основан на принципе, что аномальные точки легче изолировать: случайное разбиение пространства признаков требует fewer разбиений для изоляции аномалий.

\textbf{Рекомендации по кормлению.} Модель XGBoost рекомендует объём корма на основе надоев, \acrshort{dim}, номера лактации и текущего потребления. Базовая формула для расчёта сухого вещества (СВ):

$$\text{СВ} = \text{Базовая\_потребность} + 0.45 \cdot \text{Надой\_кг}$$

Модель корректирует базовую формулу с учётом номера лактации (первотелки требуют больше энергии), \acrshort{dim} (пик лактации требует максимального потребления) и текущего \acrshort{bcs}.

\subsubsection{Инженерия признаков}

Для каждой модели формируется набор признаков (features) на основе данных из PostgreSQL. Признаки включают лаговые (надои за предыдущие 7, 14, 30 дней), скользящие статистики (среднее, стандартное отклонение, минимум, максимум за окна 7, 14, 30 дней), производные метрики (\acrshort{fpr}, абсолютное и процентное изменение надоев, \acrshort{dim}, номер лактации) и сезонные признаки (месяц, день недели).

Обработка пропущенных значений выполняется стратегией заполнения, зависящей от типа признака: медиана для числовых, мода для категориальных, ноль для счётчиков. Для лаговых признаков у новых животных (с недостаточной историей) пропущенные значения заполняются медианой по стаду.

\subsubsection{MLOps: обучение, переобучение и мониторинг}

Автоматизация жизненного цикла моделей охватывает несколько аспектов. Автообучение запускается при отсутствии файла модели при запуске сервиса; минимальное количество записей для обучения --- 100. Плановое переобучение выполняется через APScheduler с настраиваемым расписанием (по умолчанию еженедельно, в воскресенье в 02:00). Оптимизация гиперпараметров осуществляется библиотекой Optuna~\cite{Optuna} с TPE-сэмплером (Tree-structured Parzen Estimator) в 8-мерном пространстве (learning\_rate, max\_depth, n\_estimators, min\_child\_weight, subsample, colsample\_bytree, reg\_alpha, reg\_lambda). Мониторинг дрейфа отслеживает prediction drift (z-среднее прогнозов > 2.0) и feature drift (z-среднее признаков > 2.5); при обнаружении дрейфа модель помечается для скорого переобучения. Версионирование обеспечивает хранение версий моделей в файловой системе с возможностью отката. Объяснимость прогнозов реализована через SHAP~\cite{SHAP} (SHapley Additive exPlanations), показывающий вклад каждого признака. Отслеживание экспериментов осуществляется через интеграцию с MLflow~\cite{MLflow} для логирования параметров обучения, метрик качества и артефактов моделей.

\subsection{Предиктивная аналитика в серверной части}

Помимо модуля машинного обучения, серверная часть содержит собственные алгоритмы аналитики, реализованные на Rust.

\textbf{Метод Хольта--Винтерса}~\cite{HoltWinters} --- тройное экспоненциальное сглаживание для прогнозирования временных рядов надоев. Модель учитывает три компоненты: уровень (level), тренд (trend) и сезонность (seasonality). Рекуррентные формулы:

\begin{align}
l_t &= \alpha \cdot (y_t - s_{t-m}) + (1 - \alpha) \cdot (l_{t-1} + b_{t-1}) \\
b_t &= \beta \cdot (l_t - l_{t-1}) + (1 - \beta) \cdot b_{t-1} \\
s_t &= \gamma \cdot (y_t - l_t) + (1 - \gamma) \cdot s_{t-m}
\end{align}

где $l_t$ --- уровень, $b_t$ --- тренд, $s_t$ --- сезонная компонента, $m$ --- период сезонности (7 дней для надоя), $\alpha$, $\beta$, $\gamma$ --- параметры сглаживания. Прогноз на $h$ шагов вперёд:

$$\hat{y}_{t+h} = l_t + h \cdot b_t + s_{t+h-m}$$

Для обнаружения выбросов используется межквартильный размах (IQR): значение считается выбросом, если оно выходит за пределы $[Q_1 - 1.5 \cdot \text{IQR}, Q_3 + 1.5 \cdot \text{IQR}]$. Для обнаружения структурных сдвигов применяется t-критерий: если среднее последних 7 дней значимо отличается от среднего предыдущих 14 дней (p-value < 0.05), фиксируется структурный сдвиг.

\textbf{Кривая лактации Вуда}~\cite{LactationWood} --- лог-линейная модель, описывающая зависимость надоя от дня лактации:

$$y(t) = a \cdot t^b \cdot e^{-ct}$$

Логарифмирование приводит к линейному виду:

$$\ln y(t) = \ln a + b \cdot \ln t - c \cdot t$$

Параметры $\ln a$, $b$ и $c$ оцениваются методом наименьших квадратов. Модель позволяет прогнозировать надои на основе дня лактации, определять день пика лактации $t_{\max} = -b/c$ и сравнивать фактические надои с ожидаемыми для выявления отклонений.

\textbf{Индекс здоровья} --- составной показатель на основе z-оценок:

$$\text{Health Index} = -\frac{1}{n} \sum_{i=1}^{n} \left|z_i\right|$$

где $z_i = (x_i - \mu_i) / \sigma_i$ --- стандартизованное отклонение $i$-го показателя (надои, \acrshort{scc}, активность, жвачка) от среднего по стаду. Индекс принимает значения от $-\infty$ до 0, при этом значения ниже $-2$ указывают на потенциальные проблемы со здоровьем.

\textbf{Оптимизатор запуска} --- определяет оптимальную дату сухостойного периода на основе прогноза лактации и планового отёла. Стандартный период запуска --- 60 дней до отёла. Оптимизатор корректирует дату с учётом текущего надоя: если надой выше порога (например, 10 кг/день), запуск может быть сдвинут на более позднюю дату.

\subsubsection{Сравнение моделей прогнозирования временных рядов}

Для повышения точности прогнозирования надоев в серверной части реализована система автоматического сравнения 8~моделей прогнозирования временных рядов. Система позволяет для каждой коровы выбрать модель, лучше всего соответствующую её индивидуальной динамике надоев.

Исторические данные (90~дней) разбиваются на обучающую (80~процентов) и тестовую (20~процентов) выборки хронологически. Каждая модель обучается на обучающей выборке, после чего вычисляются метрики качества прогноза на тестовой выборке:

\textbf{MAPE} (Mean Absolute Percentage Error) --- средняя абсолютная процентная ошибка:
$$\text{MAPE} = \frac{100\%}{n} \sum_{i=1}^{n} \frac{|y_i - \hat{y}_i|}{|y_i|}$$

Значение MAPE менее 10~процентов означает отличное качество прогноза, 10--20~процентов --- хорошее, более 20~процентов --- неудовлетворительное. MAPE выбран как основная метрика, поскольку она инвариантна к масштабу (надои разных коров могут отличаться в 2--3 раза).

\textbf{RMSE} (Root Mean Squared Error) --- среднеквадратичная ошибка (в литрах):
$$\text{RMSE} = \sqrt{\frac{1}{n} \sum_{i=1}^{n} (y_i - \hat{y}_i)^2}$$

RMSE дополняет MAPE, показывая типичное отклонение прогноза в абсолютных единицах. Большее, чем MAPE, штрафование крупных ошибок делает RMSE чувствительным к выбросам.

Лучшей моделью признаётся та, которая демонстрирует минимальный MAPE на тестовой выборке. Данный подход основан на предположении, что модель, лучше работающая на недавних данных, будет лучше прогнозировать и будущее.

Реализованы следующие модели:

\textbf{1. SES} (Simple Exponential Smoothing) --- простое экспоненциальное сглаживание. Модель поддерживает только уровень (level) и не учитывает тренд или сезонность. Рекуррентная формула:
$$\hat{y}_{t+1} = \alpha \cdot y_t + (1 - \alpha) \cdot \hat{y}_t$$
где $\alpha \in (0, 1)$ --- параметр сглаживания, оптимизируемый перебором по сетке $\{0.05, 0.10, \ldots, 0.95\}$ с выбором минимального SSE (sum of squared errors). Для устранения плоского прогноза вычисляется подразумеваемый наклон (implied slope) по последним fitted-значениям и экстраполируется линейно.

\textbf{2. SMA} (Simple Moving Average) --- простое скользящее среднее с окном $w$ (по умолчанию 7~дней). Прогноз --- среднее последних $w$ наблюдений. Как и SES, дополняется implied slope для формирования наклонного прогноза.

\textbf{3. WMA} (Weighted Moving Average) --- взвешенное скользящее среднее. Веса назначаются линейно возрастающими: $w_i = i$ для $i = 1, \ldots, w$, что придаёт большую значимость недавним наблюдениям. Прогноз:
$$\hat{y}_{t+1} = \frac{\sum_{i=1}^{w} w_i \cdot y_{t-w+i}}{\sum_{i=1}^{w} w_i}$$

\textbf{4. Linear Regression} --- линейная регрессия по времени. Модель оценивает параметры $\hat{y}_t = a + b \cdot t$ методом наименьших квадратов. Подходит для данных с выраженным линейным трендом, но не улавливает сезонность.

\textbf{5. Holt's Linear Method} --- двойное экспоненциальное сглаживание, добавляющее компоненту тренда к SES:
\begin{align}
l_t &= \alpha \cdot y_t + (1 - \alpha) \cdot (l_{t-1} + b_{t-1}) \\
b_t &= \beta \cdot (l_t - l_{t-1}) + (1 - \beta) \cdot b_{t-1}
\end{align}
Прогноз: $\hat{y}_{t+h} = l_t + h \cdot b_t$. Параметры $\alpha$, $\beta$ оптимизируются jointly по сетке $10 \times 10$ значений.

\textbf{6. Holt--Winters} --- тройное экспоненциальное сглаживание (см.~формулы~(1)--(3) выше). Добавляет сезонную компоненту с периодом $m = 7$~дней (недельная сезонность надоев). Наиболее гибкая модель, способная улавливать уровень, тренд и сезонность одновременно.

\textbf{7. Fourier Series} --- модель на основе рядов Фурье с линейным трендом. Данные сначала детрендируются (линейная регрессия), затем остатки аппроксимируются суммой синусов и косинусов с фиксированным базовым периодом 7~дней (недельная сезонность):
$$\hat{y}_t = (a + b \cdot t) + \sum_{k=1}^{K} \left[ A_k \cos\!\left(\frac{2\pi t}{T_k}\right) + B_k \sin\!\left(\frac{2\pi t}{T_k}\right) \right]$$
где $T_k = 7 / k$ --- период $k$-й гармоники, $K$ --- количество гармоник (по умолчанию 3). Коэффициенты $A_k$, $B_k$ оцениваются методом наименьших квадратов. Модель хорошо улавливает периодические паттерны, но может экстраполировать тренд слишком агрессивно.

\textbf{8. ARIMA(0,1,1)} --- модель авторегрессии интегрированного скользящего среднего. Данная спецификация эквивалентна экспоненциальному сглаживанию с дрифтом. Модель работает с первыми разностями $\Delta y_t = y_t - y_{t-1}$:
$$\Delta \hat{y}_t = \mu + \theta \cdot \varepsilon_{t-1}$$
где $\mu$ --- средний дрифт (среднее приращение), $\theta$ --- параметр MA(1), $\varepsilon_{t-1}$ --- предыдущая ошибка. Параметр $\theta$ оптимизируется перебором по сетке $\{0.00, 0.05, \ldots, 0.95\}$. Для прогноза на $h \geq 2$ шагов вклад MA-члена равен нулю, и прогноз вырождается в прямую с наклоном $\mu$.

Интерфейс системы отображает сравнительную таблицу моделей с их MAPE и RMSE, а также интерактивный график, позволяющий пользователю переключаться между моделями и визуально оценивать качество подгонки (fitted values) и прогноза.

\subsubsection{Ансамблевый прогноз}

Для повышения робастности прогнозирования реализован ансамблевый эндпоинт \texttt{/analytics/ensemble-forecast}, объединяющий прогноз ML-сервиса (XGBoost) и лучшую статистическую модель (Rust) методом взвешенного среднего. Вес каждой модели определяется обратно пропорционально её MAPE:

$$w_{\text{ML}} = \frac{1/\text{MAPE}_{\text{ML}}}{1/\text{MAPE}_{\text{ML}} + 1/\text{MAPE}_{\text{Rust}}}, \quad w_{\text{Rust}} = 1 - w_{\text{ML}}$$

Итоговый прогноз для дня $t$:

$$\hat{y}_t = w_{\text{ML}} \cdot \hat{y}_t^{\text{ML}} + w_{\text{Rust}} \cdot \hat{y}_t^{\text{Rust}}$$

ML-сервис имеет оценочный MAPE 8~процентов (настройка по умолчанию), тогда как MAPE Rust-модели вычисляется динамически для каждого животного. Если ML-сервис недоступен, система автоматически откатывается к прогнозу лучшей Rust-модели с весом 1.0. Интерфейс отображает веса обеих моделей, а также вклад каждой в итоговый прогноз, что повышает прозрачность и доверие фермера к системе.

\subsubsection{Объяснимость прогнозов (SHAP)}

Для обеспечения интерпретируемости прогнозов ML-сервис использует метод SHAP~\cite{SHAP} (SHapley Additive exPlanations), основанный на теории кооперативных игр Шепли. Для каждого прогноза надоя вычисляются SHAP-значения, показывающие вклад каждого признака в отклонение от базового значения модели:

$$\phi_i = \sum_{S \subseteq F \setminus \{i\}} \frac{|S|!(|F|-|S|-1)!}{|F|!} \left[ f(S \cup \{i\}) - f(S) \right]$$

где $\phi_i$ --- SHAP-значение признака $i$, $F$ --- множество всех признаков, $f(S)$ --- прогноз модели на подмножестве признаков $S$. Для XGBoost используется оптимизированный TreeExplainer, вычисляющий точные SHAP-значения за линейное время.

SHAP-объяснения включены в ответы пяти моделей прогнозирования (мастит, кетоз, охота, выбраковка, BCS) и модели прогноза надоев. В интерфейсе на странице животного SHAP-значения визуализируются в виде горизонтальной столбчатой диаграммы с цветовой кодировкой: зелёный для положительного вклада (увеличивает прогноз), красный для отрицательного (уменьшает). Это позволяет фермеру понять, какие именно факторы (например, высокий SCC, снижение жвачки) повлияли на конкретный прогноз.

\subsection{Мониторинг}

Стек мониторинга реализован на базе Grafana LGTM (Loki, Grafana, Tempo, Mimir/Prometheus). На рисунке~\ref{monitoring-arch} представлена архитектура стека мониторинга.

\begin{figure}
  \centering
  \includegraphics[width=0.95\textwidth,height=0.85\textheight,keepaspectratio]{monitoring.png}
  \caption{Архитектура стека мониторинга}
  \label{monitoring-arch}
\end{figure}

Prometheus~\cite{Prometheus} обеспечивает сбор и хранение метрик с 30-дневным хранением: счётчики HTTP-запросов (\texttt{http\_requests\_total}), гистограммы длительности запросов (\texttt{http\_request\_duration\_seconds}), метрики подключения к БД и Redis. Grafana~\cite{Grafana} визуализирует дашборды: преднастроенный дашборд backend отображает частоту запросов, время ответа, процент ошибок и использование памяти. Loki~\cite{Loki} агрегирует логи: Promtail собирает логи из Docker-контейнеров, парсит JSON-формат (уровень, сообщение, цель) и отправляет в Loki с хранением 168 часов и компактизацией. Tempo~\cite{Tempo} обеспечивает распределённую трассировку через OpenTelemetry OTLP/gRPC, позволяя отследить путь запроса от Nginx через Backend до PostgreSQL.

\subsection{Непрерывная интеграция и доставка (CI/CD)}

Конвейер CI/CD реализован на GitHub Actions~\cite{GHActions}. На рисунке~\ref{cicd} представлена схема конвейера непрерывной интеграции и доставки.

\begin{figure}
  \centering
  \includegraphics[width=0.95\textwidth,height=0.85\textheight,keepaspectratio]{ci-cd.png}
  \caption{Конвейер непрерывной интеграции и доставки}
  \label{cicd}
\end{figure}

Конвейер включает:

\textbf{Непрерывная интеграция} запускается при каждом push и pull request. Три параллельные задачи: для backend на Rust выполняются проверка форматирования (\texttt{cargo fmt --check}), статический анализ (\texttt{cargo clippy -D warnings}), аудит зависимостей (\texttt{cargo audit}) и запуск тестов (\texttt{cargo test}) с PostgreSQL в качестве сервисного контейнера; для frontend на Node.js~20 --- установка зависимостей (\texttt{npm ci}), аудит уязвимостей (\texttt{npm audit --production --audit-level=high}), линтинг (\texttt{npm run lint}), проверка форматирования (\texttt{npm run format:check}), проверка типов Svelte (\texttt{npm run check}) и запуск тестов Vitest; для analytics-ml на Python~3.12 --- установка зависимостей и запуск тестов pytest.

\textbf{Непрерывная доставка} запускается при push в main/master или при создании тега \texttt{v*}. Три параллельные задачи сборки Docker-образов с использованием Docker buildx и кэша GitHub Actions. Образы публикуются в GHCR~\cite{GHCR} с тегами \texttt{latest} и \texttt{<sha>}.

При создании тега версии дополнительно запускается задача деплоя: SSH на сервер, pull образов, создание тегов для отката, \texttt{docker compose -f docker-compose.prod.yml up -d}, проверка работоспособности (health check) в цикле: 5 попыток с интервалом 5 секунд.

\subsection{Развёртывание}

Для развёртывания используется Docker Compose с двумя конфигурациями. На рисунке~\ref{deploy} представлена схема развёртывания системы.

\begin{figure}
  \centering
  \includegraphics[width=0.95\textwidth,height=0.85\textheight,keepaspectratio]{deployment.png}
  \caption{Схема развёртывания системы}
  \label{deploy}
\end{figure}

\textbf{Разработческая} (\texttt{docker-compose.yml}): 7 сервисов --- PostgreSQL (TimescaleDB), Redis, backend, frontend, Nginx, analytics-ml, MLflow. Nginx на порту 8080. PostgreSQL инициализируется скриптом, создающим дополнительные базы данных для ML-сервиса.

\textbf{Продуктовая} (\texttt{docker-compose.prod.yml}): те же 7 сервисов с усилением безопасности: переменные окружения через env-файлы, обязательная валидация секретов (\texttt{POSTGRES\_PASSWORD}, \texttt{JWT\_SECRET}), \texttt{restart: unless-stopped}, ротация логов (10 МБ, 3 файла), ограничения памяти (256 МБ для backend/frontend, 512 МБ для PostgreSQL/analytics-ml), образы из GHCR, SSL-терминация через Nginx с монтированием сертификатов.

Nginx в продуктовой конфигурации обеспечивает: редирект HTTP $\to$ HTTPS, SSL (TLS 1.2/1.3), gzip-сжатие, заголовки безопасности (X-Frame-Options: DENY, HSTS, X-XSS-Protection), rate limiting (30 запросов/сек для API, 5 запросов/сек для auth), кэширование статических ресурсов (30 дней для immutable-ассетов).
